% Pursuing Stacks (cotes 134-2 to 134-8): archive cote 134-2, site page 68 (PDF page 69).
% appendix (letters to Breen). Draft compared: ps1.tex:1539-1563.
% AI-transcribed by Claude Opus 5.5 (claude-opus-5-5, adaptive thinking, effort high), Message Batches API, batch ps1 (msgbatch_01Rv5hCaNTtUt4JHx1yqoHPL), 2026-11-12; not yet checked by eye.
\scanpage{134-2/68}
\mspage{-\hand{5}3-}

(for example the étale topos of a scheme). Doubtless you must be confusing
the homotopical algebra which one can perform on $X$ \over{;}{\unsure{,}} using semi-simplicial
sheaves, \deleted{the} stacks of all kinds, the relations between these \inserted{-} and the other
point of view according to which $X$ (with its very rich structure of topos)
virtually disappears so as to become no more than a pale element of a
``homotopical category'' (or pro-homotopical), deduced from the topos by a
very thorough process of \hand{``}localisation\hand{''}\over{,}{.} \over{at}{At} \deleted{the} first sight, all that still
\deleted{comes from this}\inserted{remains with} poor \inserted{stripped} $X$ , are the $\pi_i$ - and its cohomology groups with constant
coefficients - or at the worst twisted constant coefficients. When one digs
more into this definition of ``what is left to this poor $X$ '' one falls
precisely on \textit{the locally constant} n-\textit{stacks} (as an $f : X \to X'$ which is a
homotopy equivalence induces a $(n + 1)$-equivalence between the categories
of locally constant n-stacks on $X$ and on $X'$) - which of course contain
the abelian chain complexes of length $n$ of sheaves with locally constant
cohomology sheaves, and the hyper-cohomology of these. It is thus that one
arrives at this triangle of objects which mutually determine each other:

% NOTE: the three arrows of the diagram appear to be drawn by hand
\diagram{
\begin{tikzcd}[column sep=-3em]
 & \begin{tabular}{c}topos (or topological space\\ or semi-simplicial complex)\\ \textit{modulo} n-homotopy\end{tabular} \arrow[dl, leftrightarrow] \arrow[dr, leftrightarrow] & \\
\begin{tabular}{l}n-groupoids\\ (up to n-equivalence)\end{tabular} \arrow[rr, leftrightarrow] & & \begin{tabular}{l}``special''n-categories\\ (up to n-equivalence).\end{tabular}
\end{tikzcd}
}

One says that an n-category $E_n$ is ``special'' (or n-\textit{galois}) if it is
n-equivalent to the category of locally constant $(n-1)$-stacks on an appropriate
topological space (or a topos), or, what \deleted{must}\inserted{should} be equivalent, if it is
n-equivalent to the category of n-functors $G_n \to (n - \mathrm{Cat})$ , where $G_n$ is
an n-groupoid. If $X$ , $G_n$ , $E_n$ correspond in this way, one calls $G_n$ \uline{the}
\uline{fundamental n-groupoid} of $X$ , or of $\over{\unsure{F}}{E}_n$ , or says that $E_n$ is \uline{the category}
\uline{of local \deleted{n}\inserted{(n-1)}-systems} on $X$ , or on $G_n$ , or that $X$ is \uline{the geometric realisation}
\margin{/} of $G_n$ or of $E_n$ \over{,}{.} In analogy with the familiar case $n = 1$ , \deleted{$E_n$} \deleted{can be}\inserted{it should be}
\inserted{possible to} interprete\deleted{d} \inserted{$G_n$} as the full sub-n-category of $\mathit{Hom}_n(E_n , (n - \mathrm{cat}))$ formed by
the n-functors $E_n \to (n\text{-}\mathrm{cat})$ satisfying certain exactness properties (one

\margin{68}
